\uuid{f7hD}
\exo7id{4470}
\titre{exo7 4470}
\auteur{quercia}
\organisation{exo7}
\datecreate{2010-03-14}
\isIndication{false}
\isCorrection{true}
\chapitre{Série numérique}
\sousChapitre{Autre}
\cpge{mpsi-mp2i.procedes-sommatoires-discrets.a, pcsi.series-numeriques.a, ptsi.series-numeriques.a, mpsi-mp2i.nombres-reels-suites.d, pcsi.nombres-reels-suites.c, ptsi.nombres-reels-suites.c, bcpst1.suites-reelles, mp-mpi.series-numeriques-vectorielles.a, pt.series-numeriques.b, bcpst2.series-reelles}
\module{Analyse}
\niveau{L1}
\difficulte{}

\contenu{
\texte{
Soit $\sum u_n$ une série convergente.
On pose $v_n = \frac{u_n}{1} + \frac{u_{n-1}}{2} + \dots + \frac{u_0}{2^n}$.
}
\begin{enumerate}
    \item \question{Montrer que $v_n \to 0$ lorsque $n\to\infty$.}
\reponse{Césaro.}
    \item \question{Montrer que $\sum v_n$ converge et donner sa valeur.}
\reponse{$v_0 + v_1 + \dots + v_n = 2(u_0 + u_1 + \dots + u_n) - v_n$.}
\end{enumerate}
}
